\section{从哲学到 AI 对齐}

Amanda 在牛津和纽约大学受训为哲学家，博士研究``包含无限多人的世界中的伦理学''——这是技术性/理论性伦理学，而非应用伦理学。哲学训练让她能够自如地穿越不同领域（数学、化学、政治、伦理），这种跨领域思维能力在 AI 对齐工作中极为宝贵。

她从学术界转向工业界的动机很简单：\textit{``如果你尝试做有影响力的事而失败了，你至少尝试过了，然后可以回去做学者。''} 2017-2018 年 AI 正在成为大事件但尚未被广泛认知，Amanda 认为这是一个理想的切入时机。

她的职业路径经历了三个阶段：AI 政策（政治影响）$\to$ AI 评估（比较模型与人类输出）$\to$ 技术对齐（加入 Anthropic）。她发现自己在技术领域比在政策领域更如鱼得水：\textit{``政治是混乱的，更难找到确定性的、清晰的、可证明的、美丽的解决方案。''}

\begin{knowledgebox}{打破``技术 vs 非技术''二元对立}
Amanda 强调人们错误地创造了``技术人''和``非技术人''的身份二元对立。\textit{``我认为很多人其实完全有能力做这类工作，只要他们愿意尝试。''} 她的建议：找一个项目，动手做——模型现在已经强大到可以辅助学习，技术门槛比以前低得多。
\end{knowledgebox}

\subsection{Claude 的人格：对齐项目而非产品需求}

Amanda 与 Claude 的对话次数比 Anthropic 任何人都多——用 5-6 种不同方法与 Claude 交互，不仅仅是著名的 Slack 频道。

人格设计被构思为\textbf{对齐}项目而非产品考量。核心问题：\textit{``想象你把一个人放在这个位置——他们知道自己将与数百万人交谈……你希望他们以一种非常丰富的意义上表现良好。''}

\begin{importantbox}{世界旅行者框架}
\textit{``想象一个理想的人，能环游世界，与各种不同的人交谈，几乎每个人都会说'这真是一个很好的人'。''} 这个人：
\begin{itemize}
    \item 真诚——不伪装
    \item 坦率表达意见——但不强加
    \item 开放——愿意理解不同观点
    \item 尊重他人自主权——不试图改变别人
\end{itemize}

关键：这个人\textbf{不}会入乡随俗地改变价值观——那反而是\textit{``某种程度的不尊重''}。Claude 的性格应该有一致性，不因对话者不同而变。
\end{importantbox}

这根植于亚里士多德式的``好品格''理念——不仅是伦理规则的遵守，还包括细微、幽默、关怀和尊重自主权。Amanda 区分了``丰富的品格''（rich character）和``薄弱的伦理''（thin ethics）：后者只有规则，前者有判断力。

\subsection{Claude 与政治/争议话题}

Amanda 认为价值观更像物理学而非味觉偏好——是值得探究的对象，而非固定的偏好。Claude 应该理解世界上所有的价值观，对它们保持好奇，但不谄媚。

\textit{``知识谦逊——说话的欲望会迅速下降。''} 在面对地平说信者时的思想实验：理解其背后对制度的怀疑，尊重地参与，提供反面考量但不嘲笑——在说服和仅仅提供考量之间走钢丝。

\subsection{本章小结}

Claude 的人格设计源于对齐研究而非产品需求，用``世界旅行者''框架定义理想行为者。好品格的关键是一致性和判断力，而非规则的机械遵守。

